\subsection{Verifier-Guided Reasoning (Interwhen)}
\label{sec:bg_interwhen}

Interwhen~\citep{bhat2026interwhen} is a verifier-guided reasoning
framework in which a monitor periodically polls an LLM's reasoning
trace, forks the model to extract verifiable state from the partial
trace, and runs verifiers on that state asynchronously --- injecting
feedback and re-prompting only when a violation is detected. The
verifiers are code-based predicates over partial traces and can be
synthesized automatically from natural-language policy documents,
including provably correct variants in Lean and z3. Interwhen is
evaluated across both non-agentic reasoning (math, logic) and agentic
tool-calling benchmarks (e.g., $\tau^2$-bench). Crucially for our
adaptation, the state the verifier consumes is itself recovered by an
LLM extraction step over the trace, and interwhen explicitly analyzes
the framework's robustness to errors in that extraction.

We adapt the mechanism rather than re-implement it, preserving
interwhen's verify/fix control flow exactly --- on flag, inject
feedback and re-prompt; on clean, dispatch the tool and continue.
Three things change. First, instead of polling the trace at intervals,
we fix the intercept at a single deterministic commit point: the
moment the primary model emits a \verb|<tool_call>| block, whose parsed
JSON is the state under check. Second, the LLM extraction reads the
patient \emph{case} --- the task input --- rather than the model's own
reasoning trace; interwhen recovers the model's asserted state from the
trace and checks it against a policy, whereas we recover an independent
reference from the case and check the model's proposed inputs against
it. Third, in place of a policy-synthesized verifier we use a
hand-built deterministic field-by-field comparator. The transport also
differs: we use vLLM's text-completions endpoint with a custom
synchronous loop driver rather than the streaming variant, because the
primary model is served behind a standard completions API.

\subsection{Clinical Tool-Calling and KARMA}
\label{sec:bg_karma}

KARMA~\citep{karma2024} is a structured evaluation framework for medical
AI published by EkaCare. We use its publicly released
medical\_calculator\_eval benchmark ($N=1066$ vignettes spanning 24
calculator categories, from cardiology to obstetrics) and the
associated MedAI Model Context Protocol (MCP)
toolset~\citep{medai_mcp}.

The MCP toolset exposes six meta-tools, of which three are central to
this study:
\toolname{medical\_calculator\_list} enumerates available calculators;
\toolname{medical\_calculator\_input} returns the JSON Schema for a chosen
calculator's input; and \toolname{medical\_calculator\_output} dispatches the
calculator with an \verb|input_data| payload. Each calculator carries a
typed JSON Schema with enums, ranges, and field descriptions. The union
of input field names across all calculators yields a vocabulary of
$\sim$500 unique fields. This union is the reference vocabulary the
verifier and extractor share, by construction --- both consume the
schemas EkaCare ships, so neither can refer to a field the other does
not know about.

\subsection{Verification Surfaces: Where to Intervene}
\label{sec:bg_surfaces}

LLM-driven clinical AI admits intervention at several surfaces, which
are often conflated in casual discussion. We distinguish:

\textit{Input verification} (this paper). The verifier inspects the
model's proposed tool-call arguments \emph{before} dispatch. Catches
wrong inputs at the boundary where they would otherwise propagate
into a deterministic backend. Fires at intermediate commit points;
firing frequency scales with tool-call rate.

\textit{Output verification.} The verifier reads the model's final
textual answer and may request revision. Cannot catch a wrong tool
input that produced a confidently-stated wrong answer unless the
final text exposes the inconsistency. Condition D in our study is the
closest output-style comparator.

\textit{Grounding / retrieval-augmented verification.} The reference
is retrieved text. A different problem class --- the question becomes
``is the claim supported by retrieved evidence?'' rather than ``is the
proposed input consistent with the case?'' Not evaluated here.

\textit{Self-verification.} The model is asked, in-prompt, to check
itself. No external reference. Condition C is our control for ``is any
external machinery necessary?''

This paper concerns the first surface. Every claim about ``the
verifier'' refers to procedural input verification specifically.

Procedural input verification is a \emph{special case} of
reasoning-trace verification: the model's tool-call emission is one
step in its trace, and we intercept at that step. What distinguishes
our setting is where the verifier's reference comes from and where the
commit point sits. Interwhen recovers verifiable state by polling the
trace at intervals; in a clinical agent most of that trace is
unstructured prose, and the reference the verifier needs --- the
patient's actual values --- is not in the trace at all but in the case
text. The model's typed tool-call payload is the first point in the
trace where a machine-checkable structured commit exists, and the
calculator's input schema defines exactly what fields constitute that
commit. A flag therefore
localizes to a specific (calculator, field, value) tuple --- auditable
in the way clinical safety reviews require. Second, existing clinical
decision-support pipelines already \emph{have} these tool boundaries
(risk calculators, formularies, eligibility checkers, FHIR queries),
so input verification inserts at a place the deployment architecture
already exposes, rather than requiring the workflow to be restructured
around inline reasoning checks.

\subsection{The Reference-Extraction Problem}
\label{sec:bg_reference}

A subtlety in the clinical adaptation deserves explicit attention:
\emph{the reference itself must come from somewhere}. Interwhen already
relies on an LLM to extract the verifier's input variables, but it
extracts them from the model's own reasoning trace --- the state the
verifier checks is the model's asserted state, scored against a policy.
In our setting the verifier instead needs the patient's true values,
which live only in the natural-language case. Before the verifier can
compare anything, the case must be parsed into field--value pairs by an
LLM, and that extracted record is then treated as the reference against
which the model's proposed inputs are judged. The model-based
extraction step is therefore not new to interwhen; what is new is that
it now stands between the case and the verifier as the source of ground
truth, which is where the verifier's safety guarantee can weaken.

We use an external LLM of a different model family than the primary
to produce the reference (specific model, version, and serving
configuration in \cref{sec:methods_benchmark}). The cost of this
design is that the extractor is itself an LLM, with its own
distribution of errors. This paper's main finding is that those
errors --- not the verifier mechanism --- bound the achievable
accuracy.

\subsection{Which Failure Classes Input Verification Can Reach}
\label{sec:bg_scope}

EkaCare's own error analysis of
medical\_calculator\_eval~\citep{ekacare_error_analysis} identifies two
failure classes that persist even with tool access. The first is
\emph{wrong calculator selected} --- the model picks a semantically
adjacent but distinct tool (e.g., a generic BMI calculator where a
sex-specific variant is required). The second, affecting 12--13\% of
samples across the three models they evaluated, is \emph{right
calculator, wrong inputs}, which decomposes into (a) enum misselection
(mapping ``sedentary'' to ``very\_active'' in a TDEE activity field)
and (b) unit-conversion overreach, where the model receives a correct
calculator result and then rewrites it into a different unit,
introducing an error the calculator did not make.

Procedural input verification is scoped to exactly one of these
sub-classes. It fires before dispatch on the tool-call payload, so it
can catch enum misselection and analogous wrong-input errors at the
boundary --- sub-class (a) above. It \emph{cannot} reach sub-class (b):
unit-conversion overreach occurs after the calculator returns, on the
output surface, where only a post-hoc check such as Condition D could
intervene. Wrong-calculator selection is likewise out of scope --- the
verifier checks field values for the calculator the model has already
chosen, not the choice of calculator itself. The addressable ceiling
for the primary study conditions is therefore the enum/wrong-input
portion of the 12--13\% band, not the full band; we report the split
where the logs permit and calibrate expected effect sizes accordingly.
Condition D is the only arm positioned to reach the unit-overreach
sub-class, and it targets a different failure than the B$'$+E
conditions rather than a more thorough version of the same one.
